\documentclass[11pt,a4paper]{report}

% ============================================================================
% RetailVision Edge AI — Manual Técnico
% Universidad Santo Tomás, Seccional Tunja
% ============================================================================

\usepackage[utf8]{inputenc}
\usepackage[T1]{fontenc}
\usepackage{lmodern}
\usepackage[spanish,es-noindentfirst]{babel}
\usepackage[a4paper,margin=2.5cm]{geometry}
\usepackage{xcolor}
\usepackage{graphicx}
\usepackage{tabularx}
\usepackage{longtable}
\usepackage{booktabs}
\usepackage{listings}
\usepackage{enumitem}
\usepackage{fancyhdr}
\usepackage{titlesec}
\usepackage{hyperref}
\usepackage{caption}
\usepackage{array}
\usepackage{colortbl}
\usepackage{microtype}
\sloppy

% --- Paleta (coherente con el dashboard: acento ámbar/naranja) ------------
\definecolor{rvAccent}{HTML}{C76A00}
\definecolor{rvAccentDark}{HTML}{8A4600}
\definecolor{rvInk}{HTML}{2C3242}
\definecolor{rvMuted}{HTML}{5B6478}
\definecolor{rvLine}{HTML}{CCD4E2}
\definecolor{rvBg}{HTML}{F7F8FA}
\definecolor{rvCodeBg}{HTML}{F2F3F6}

% --- Tipografía de títulos --------------------------------------------------
\titleformat{\chapter}[display]
  {\normalfont\bfseries\color{rvInk}}
  {\color{rvAccent}\Large\chaptertitlename\ \thechapter}
  {6pt}{\Huge}
\titleformat{\section}
  {\normalfont\Large\bfseries\color{rvInk}}{\thesection}{0.6em}{}
\titleformat{\subsection}
  {\normalfont\large\bfseries\color{rvAccentDark}}{\thesubsection}{0.6em}{}

% --- Encabezados -------------------------------------------------------------
\setlength{\headheight}{14.5pt}
\addtolength{\topmargin}{-0.5pt}
\pagestyle{fancy}
\fancyhf{}
\renewcommand{\headrulewidth}{0.4pt}
\renewcommand{\footrulewidth}{0pt}
\fancyhead[L]{\footnotesize\color{rvMuted} RetailVision Edge AI}
\fancyhead[R]{\footnotesize\color{rvMuted} \leftmark}
\renewcommand{\chaptermark}[1]{\markboth{#1}{}}
\fancyfoot[C]{\small\color{rvMuted} \thepage}

% --- Bloques de código -------------------------------------------------------
\lstdefinestyle{rvcode}{
  backgroundcolor=\color{rvCodeBg},
  basicstyle=\ttfamily\footnotesize,
  keywordstyle=\color{rvAccentDark}\bfseries,
  commentstyle=\color{rvMuted}\itshape,
  stringstyle=\color{rvAccent},
  numberstyle=\tiny\color{rvMuted},
  numbers=left,
  numbersep=8pt,
  frame=single,
  rulecolor=\color{rvLine},
  breaklines=true,
  breakatwhitespace=false,
  showstringspaces=false,
  tabsize=2,
  xleftmargin=12pt,
  xrightmargin=4pt,
  framexleftmargin=8pt,
}
\lstset{style=rvcode}

\hypersetup{
  colorlinks=true,
  linkcolor=rvAccentDark,
  urlcolor=rvAccentDark,
  citecolor=rvAccentDark,
  pdftitle={RetailVision Edge AI -- Manual Tecnico},
  pdfauthor={Equipo RetailVision Edge AI -- Universidad Santo Tomas, Seccional Tunja},
}

\newcommand{\rvnota}[1]{%
  \par\noindent\fcolorbox{rvAccent}{rvCodeBg}{%
    \parbox{0.95\linewidth}{\small \textbf{\color{rvAccentDark}Nota:} #1}%
  }\par\vspace{4pt}
}

\setlist[itemize]{leftmargin=1.4em}
\setlist[enumerate]{leftmargin=1.6em}

\title{\textbf{\Huge RetailVision Edge AI}\\[6pt]
  \Large Manual Técnico\\[4pt]
  \large Sistema de Visión Computacional para \textit{Space Management}\\
  y Analítica de Tráfico en Retail}
\author{Equipo RetailVision Edge AI\\
  Universidad Santo Tomás, Seccional Tunja\\
  Proyecto de Investigación / Semillero\\
  Docente asesor: Sergio Arley Puerto Moreno}
\date{2026}

\begin{document}

\begin{titlepage}
\centering
\vspace*{1.5cm}
{\Huge\bfseries\color{rvAccent} RetailVision Edge AI\par}
\vspace{0.4cm}
{\Large\color{rvInk} Manual Técnico\par}
\vspace{0.3cm}
{\large\color{rvMuted} Sistema de Visión Computacional para \textit{Space Management}\\
y Analítica de Tráfico en Retail\par}
\vspace{2.2cm}

{\large\color{rvInk}\textbf{Universidad Santo Tomás, Seccional Tunja}\par}
\vspace{0.2cm}
{\normalsize\color{rvMuted} Proyecto de Investigación / Semillero\par}
\vspace{0.15cm}
{\normalsize\color{rvMuted} Docente asesor: Sergio Arley Puerto Moreno\par}

\vspace{2.2cm}
{\normalsize\color{rvInk}\textbf{Equipo}\par}
\vspace{0.3cm}
\begin{tabular}{l}
Silvana Sofía Siza Soriano \\
Yeikov Jared Ramos Sánchez \\
Sarita Toloza Buitrago \\
Juan David Sandoval Bello \\
Thomas Felipe Cepeda Mogollón \\
Wendy Mariana Vargas Ramírez \\
Andrés Felipe Galeano Téllez \\
Sara Valentina Caicedo Duarte \\
\end{tabular}

\vfill
{\large\color{rvMuted} 2026\par}
\end{titlepage}

\tableofcontents
\clearpage

% ============================================================================
\chapter{Introducción}
% ============================================================================

\section{Propósito del sistema}

\textbf{RetailVision Edge AI} es un sistema de visión computacional que transforma el video de
cámaras de seguridad de una tienda en decisiones automatizadas de \textit{space management}.
A partir del flujo de video, el sistema:

\begin{itemize}
  \item Detecta y rastrea personas frente a las góndolas (pose + tracking).
  \item Identifica interacciones de \textbf{toma} (\textit{pick-up}) y \textbf{devolución}
        (\textit{put-back}) de producto.
  \item Mide tiempos de permanencia y tráfico por zona (heatmap).
  \item Clasifica cada zona como \textit{caliente}, \textit{templada} o \textit{fría} según su
        actividad observada.
  \item Genera recomendaciones de planograma mediante un agente determinista (reglas sobre datos
        observados, sin modelos de lenguaje ni llamadas a servicios externos).
\end{itemize}

Todo el procesamiento ocurre \textbf{100\,\% en el dispositivo local} (\textit{edge}), bajo un
enfoque de \textbf{privacidad por diseño}: los rostros se anonimizan dentro del propio pipeline,
antes de que cualquier fotograma se renderice o se persista, en cumplimiento de la Ley 1581 de
2012 (Colombia) de protección de datos personales.

\section{Contexto académico}

El proyecto se desarrolla en el marco de un semillero de investigación de la \textbf{Universidad
Santo Tomás, Seccional Tunja}, bajo la asesoría del docente \textbf{Sergio Arley Puerto Moreno}.
Su primera versión se presentó como proyecto de la hackathon \textit{HackTech 5.0} bajo el nombre
``Góndola Inteligente''; este manual documenta la versión consolidada del sistema, ya en proceso
de inscripción al Registro de Software Académico de la institución.

\section{Alcance de este documento}

Este manual está dirigido a \textbf{desarrolladores, mantenedores y evaluadores técnicos} del
proyecto. Describe la arquitectura de software, el stack tecnológico, el procedimiento de
instalación y ejecución, el contrato de la API REST, el modelo de datos exportado, la lógica de
dominio (máquina de estados de interacción), la estrategia de privacidad, y la suite de pruebas y
calidad. Para una guía orientada al uso operativo del dashboard (sin detalles de implementación),
véase el \textbf{Manual de Usuario} que acompaña este documento.

% ============================================================================
\chapter{Arquitectura del sistema}
% ============================================================================

\section{Visión general}

El sistema sigue una \textbf{arquitectura por capas (Clean Architecture / Hexagonal ligera)}: el
dominio contiene la lógica de negocio pura y no depende de ningún framework; los adaptadores de
infraestructura aíslan YOLO, OpenCV y el transporte HTTP. La regla de dependencia se cumple de
adentro hacia afuera: \texttt{domain} no importa nada de \texttt{infrastructure} ni
\texttt{interfaces}.

\begin{figure}[h]
\centering
\begin{lstlisting}[language=,basicstyle=\ttfamily\footnotesize,frame=single,backgroundcolor=\color{rvCodeBg},numbers=none,xleftmargin=0pt,framexleftmargin=6pt]
video de tienda
  -> YOLO-pose + ByteTrack + YOLO producto-en-mano   [infrastructure]
  -> maquina de estados (pick-up / put-back)         [domain]
  -> heatmap de trafico + analisis por zona          [domain]
  -> recomendaciones de planograma (determinista)    [domain]
  -> JSON de metricas (metricas_retail.json)         [use_cases]
  -> backend FastAPI /api/*                          [interfaces]
  -> dashboard React (5 vistas + editor de zonas)    [frontend]
\end{lstlisting}
\caption{Flujo de datos de extremo a extremo.}
\end{figure}

\section{Capas y responsabilidades}

\begin{longtable}{>{\raggedright\footnotesize\arraybackslash}p{2.2cm} >{\raggedright\footnotesize\arraybackslash}p{6.0cm} >{\raggedright\footnotesize\arraybackslash}p{6.4cm}}
\toprule
\rowcolor{rvCodeBg}\textbf{Capa} & \textbf{Contiene} & \textbf{Depende de} \\
\midrule
\endhead
\texttt{domain/} & Lógica de negocio pura: máquina de estados del cliente, geometría de
polígonos, heatmap, agente de recomendación, puertos (\texttt{Protocol}) de detección y de
repositorio. & Nada interno. \\
\texttt{application/} & Persistencia de zonas, orquestación del ciclo de vida de un análisis
(\texttt{AnalysisOrchestrator}), rutas del proyecto. & \texttt{domain/} \\
\texttt{use\_cases/} & Orquestación del pipeline completo (frames, tracking, anonimización,
exportación de métricas). & \texttt{domain/}, \texttt{infrastructure/}, \texttt{application/} \\
\texttt{infrastructure/} & Adaptadores concretos: YOLO, ByteTrack, captura de video, blur de
privacidad, ejecución del pipeline en subproceso, repositorio JSON de zonas. & Librerías
externas únicamente. \\
\texttt{interfaces/} & FastAPI (API HTTP), CLI (argparse), validación de precisión contra
\textit{ground truth}. & Todas las anteriores. \\
\bottomrule
\end{longtable}

\section{Estructura del repositorio}

\begin{lstlisting}[language=,basicstyle=\ttfamily\footnotesize,frame=single,backgroundcolor=\color{rvCodeBg}]
app/
|-- config.py                    -> constantes de ajuste (umbrales, zonas, colores)
|-- application/
|   |-- analysis_orchestrator.py -> ciclo de vida de un analisis (estado en memoria)
|   |-- zona_repository.py       -> puerto ZonaRepository
|   |-- paths.py                 -> rutas resueltas a la raiz del repo
|   `-- zonas.py                 -> persistencia/validacion/escalado de zonas
|-- domain/
|   |-- persona.py               -> maquina de estados de la persona y de la mano
|   |-- analisis.py              -> heatmap de trafico + analizador de zonas
|   |-- agente.py                -> recomendaciones de planograma (determinista)
|   |-- constantes.py            -> vocabulario de dominio (acciones, colores, clasificacion)
|   |-- detector_personas.py     -> puerto DetectorDePersonas (Protocol)
|   |-- detector_objetos.py      -> puerto DetectorDeObjetos (Protocol)
|   |-- geometria.py             -> helpers geometricos puros
|   `-- poligonos.py             -> validacion de poligonos (ray-casting)
|-- use_cases/
|   |-- procesar_video.py        -> orquesta el pipeline de deteccion
|   `-- exportar_metricas.py     -> vuelca las metricas a JSON
|-- infrastructure/
|   |-- vision.py                -> adaptadores YOLO + captura de video
|   |-- anonimizacion.py         -> mosaico de rostros (privacidad)
|   |-- dibujo.py                -> etiquetas sobre el fotograma
|   |-- pipeline_runner.py       -> ejecuta el pipeline en subproceso aislado
|   `-- zona_repository_json.py  -> adaptador JSON del puerto ZonaRepository
`-- interfaces/
    |-- api.py                   -> FastAPI (endpoints del dashboard)
    |-- cli.py                   -> pipeline por linea de comandos
    `-- validar_embudo.py        -> validacion P/R/F1 contra ground truth

frontend/src/
|-- domain/        -> logica pura (merge de zonas, tasas, agente de recomendacion)
|-- application/    -> cliente HTTP hacia el backend
`-- presentation/   -> componentes React (dashboard, editor de zonas, graficas)
\end{lstlisting}

\section{Puertos y adaptadores (Hexagonal)}

Tres puntos de extensión concretos ilustran el patrón de puertos y adaptadores aplicado en el
proyecto:

\begin{itemize}
  \item \textbf{\texttt{ZonaRepository}} (puerto, \texttt{application/zona\_repository.py}):
    contrato de persistencia de zonas (\texttt{obtener}, \texttt{guardar}, \texttt{restablecer},
    \texttt{fuente}). Su adaptador, \texttt{JsonZonaRepository}
    (\texttt{infrastructure/zona\_repository\_json.py}), persiste en dos archivos JSON: uno de
    fábrica de solo lectura y uno de \textit{runtime} editable desde la web.
  \item \textbf{\texttt{AnalysisOrchestrator}} (\texttt{application/analysis\_orchestrator.py}):
    encapsula el estado del análisis en curso (antes variables globales del módulo de la API),
    inyectado en los \textit{endpoints} vía \texttt{Depends()} de FastAPI. Recibe como
    dependencia un \texttt{ejecutor} (puerto) que desacopla la orquestación de \textbf{cómo} se
    corre el pipeline realmente.
  \item \textbf{\texttt{DetectorDePersonas} / \texttt{DetectorDeObjetos}} (\texttt{Protocol} en
    \texttt{domain/}): contratos de tipado estructural que \texttt{DetectorPoseYObjetos}
    (\texttt{infrastructure/vision.py}) satisface sin herencia explícita. Permiten a los casos de
    uso declarar qué necesitan sin acoplarse a la clase concreta de YOLO.
\end{itemize}

\rvnota{Esta separación permite sustituir cualquier adaptador (por ejemplo, un repositorio de
zonas en base de datos en vez de JSON, o un detector distinto a YOLO) sin tocar el dominio ni los
casos de uso.}

% ============================================================================
\chapter{Stack tecnológico}
% ============================================================================

\begin{longtable}{>{\raggedright\arraybackslash}p{3cm} >{\raggedright\arraybackslash}p{11.6cm}}
\toprule
\rowcolor{rvCodeBg}\textbf{Capa} & \textbf{Tecnologías} \\
\midrule
\endhead
Visión & \texttt{ultralytics} (YOLO-pose + detector de producto-en-mano \textit{fine-tuneado}),
OpenCV, NumPy, ByteTrack (tracking multi-objeto). \\
Backend & Python 3.12, FastAPI, uvicorn, \texttt{Depends()} para inyección de dependencias. \\
Frontend & React 19, Vite 8, ECharts (gráficas interactivas), three.js (plano 3D). \\
Calidad & \texttt{unittest} / \texttt{pytest}, \texttt{scripts\_qa} (privacidad, métricas,
rendimiento), ESLint, \texttt{node:test}. \\
Empaquetado & \texttt{requirements.txt} (Python), \texttt{npm} (frontend), \texttt{arrancar.ps1}
(orquestación local en Windows). \\
\bottomrule
\end{longtable}

% ============================================================================
\chapter{Instalación y configuración}
% ============================================================================

\section{Requisitos previos}

\begin{itemize}
  \item Python 3.12 (recomendado; probado también en 3.13).
  \item Node.js 18 o superior y \texttt{npm}.
  \item (Opcional) GPU NVIDIA para aceleración CUDA --- probado en RTX 2060.
\end{itemize}

\section{Backend}

\begin{lstlisting}[language=bash]
python -m venv .venv
.venv\Scripts\activate          # Windows
pip install -r requirements.txt
\end{lstlisting}

Los modelos YOLO (\texttt{.pt}) se descargan automáticamente con \texttt{ultralytics} la primera
vez que se ejecutan (\texttt{models/yolo11n-pose.pt}); el detector de producto-en-mano
\textit{fine-tuneado} se entrena sobre el dataset de Roboflow \texttt{theft-full-2500}.

Aceleración por GPU (opcional, NVIDIA con soporte CUDA 12.6):

\begin{lstlisting}[language=bash]
pip install "torch==2.13.0+cu126" "torchvision==0.28.0+cu126" \
  --index-url https://download.pytorch.org/whl/cu126
pip uninstall -y opencv-python-headless
\end{lstlisting}

\rvnota{Sin GPU, el sistema funciona en CPU sin cambios de código; solo cambia el tiempo de
procesamiento (ver \S\ref{sec:rendimiento}).}

\section{Frontend}

\begin{lstlisting}[language=bash]
cd frontend
npm install
\end{lstlisting}

\section{Variables de entorno relevantes}

\begin{longtable}{>{\ttfamily\raggedright\footnotesize\arraybackslash}p{5.0cm} >{\raggedright\arraybackslash}p{9.4cm}}
\toprule
\rowcolor{rvCodeBg}\textbf{Variable} & \textbf{Efecto} \\
\midrule
\endhead
GONDOLA\_CORS\_ORIGINS & Orígenes permitidos por CORS (por defecto, los puertos de desarrollo de
Vite). \\
GONDOLA\_MAX\_VIDEO\_MB & Límite de tamaño por video subido (por defecto, 500\,MB). \\
\bottomrule
\end{longtable}

% ============================================================================
\chapter{Ejecución}
% ============================================================================

\section{Todo en uno (Windows)}

\begin{lstlisting}[language=bash]
powershell -ExecutionPolicy Bypass -File .\arrancar.ps1
# Dashboard  -> http://localhost:5173
# Backend API -> http://localhost:8000
\end{lstlisting}

\section{Por separado}

\begin{lstlisting}[language=bash]
python -m uvicorn backend.main:app --reload --port 8000   # backend
cd frontend && npm run dev                                 # frontend
\end{lstlisting}

\section{Pipeline por línea de comandos}

\begin{lstlisting}[language=bash]
python -m app.interfaces.cli                  # con ventana OpenCV en vivo
python -m app.interfaces.cli --no-ventana     # headless (solo metricas)
python -m app.interfaces.cli --no-ventana \
  --video video/mi_video.mp4 \
  --guardar-frames outputs/vision_frames \
  --estado-json outputs/analysis_status.json \
  --run-id demo-001
\end{lstlisting}

\section{Flujo de un análisis (desde el dashboard)}

\begin{enumerate}
  \item \textbf{Subir un video} (\texttt{POST /api/video}) --- el backend extrae un fotograma de
    \textbf{vista previa con rostros ya anonimizados}.
  \item \textbf{Calibrar las zonas} sobre esa vista previa (arrastrar los vértices de cada
    góndola en el editor web).
  \item \textbf{Ejecutar el análisis} (\texttt{POST /api/analizar}) --- el pipeline YOLO corre en
    un subproceso aislado, en segundo plano.
  \item El frontend consulta \texttt{GET /api/estado} para el progreso y \texttt{GET /api/datos}
    para las métricas resultantes.
  \item Si el usuario \textbf{vuelve a guardar las zonas}, el análisis se \textbf{re-ejecuta
    automáticamente} sobre el último video procesado.
\end{enumerate}

% ============================================================================
\chapter{API REST}
% ============================================================================

Todos los \textit{endpoints} están prefijados por \texttt{/api} y son servidos por FastAPI
(\texttt{app/interfaces/api.py}). Las respuestas son JSON salvo donde se indique lo contrario.

\begin{longtable}{>{\raggedright\arraybackslash}p{1.3cm} >{\ttfamily\raggedright\footnotesize\arraybackslash}p{3.9cm} >{\raggedright\arraybackslash}p{8.6cm}}
\toprule
\rowcolor{rvCodeBg}\textbf{Método} & \textbf{Ruta} & \textbf{Descripción} \\
\midrule
\endhead
GET & /api/estado & Estado del análisis (\texttt{inactivo}, \texttt{iniciando},
\texttt{procesando}, \texttt{completado}, \texttt{error}) y progreso. \\
GET & /api/datos & Métricas reales del pipeline (zonas, embudo, heatmap, recomendaciones). 404 si
aún no hay datos. \\
POST & /api/analizar & Ejecuta el análisis YOLO en segundo plano (video pendiente, último
analizado, o el especificado por \texttt{ruta\_video}). \\
POST & /api/reiniciar & Reinicia las métricas (borra resultados y estado previos; conserva las
zonas calibradas). \\
POST & /api/video & Sube un video y prepara la vista previa anonimizada para calibrar (no analiza
aún). \\
GET & /api/zonas & Zonas vigentes (polígono + agarro) en el \textit{frame} de referencia
1280x720, y su origen (\texttt{runtime} o \texttt{default}). \\
PUT & /api/zonas & Persiste las zonas editadas; re-ejecuta el análisis si corresponde. \\
POST & /api/zonas/restablecer & Vuelve a las zonas de fábrica. \\
GET & /api/vision-frames & Índice de los fotogramas anotados del visor ``visión de las
góndolas''. \\
GET & / & Sirve el frontend estático (\texttt{dist/}). \\
\bottomrule
\end{longtable}

\section{Ejemplo: cuerpo de \texttt{GET /api/datos} (resumido)}

\begin{lstlisting}[language=,basicstyle=\ttfamily\footnotesize]
{
  "generado_en": "2026-10-24T15:02:11",
  "video": {"fps": 24.4, "ancho": 1280, "alto": 720},
  "total_personas_registradas": 12,
  "embudo_interaccion": {
    "total_pick_ups": 9, "total_put_backs": 2, "tasa_rechazo_general": 0.22
  },
  "analisis_zonas": [
    {"id": "Estante_Izquierdo", "visitas": 14, "pick_ups": 5,
     "put_backs": 1, "clasificacion": "caliente", "dwell_promedio_seg": 6.4}
  ],
  "heatmap_global": [[0, 12, 48, "..."]],
  "recomendaciones": {
    "motor": "agente determinista (reglas sobre datos observados por YOLO)",
    "por_zona": {"Estante_Izquierdo": "..."},
    "planograma": {"accion": "reubicar", "origen": "...", "destino": "..."}
  },
  "detalle_por_persona": [
    {"id_persona": 3, "permanencia_estante_seg": 8.2, "producto_adquirido": true,
     "eventos": [{"tipo": "Producto Adquirido", "frame": 142, "zona": "Estante_Central"}],
     "trayectoria": [{"segundo": 4.1, "x": 640, "y": 300}]}
  ]
}
\end{lstlisting}

% ============================================================================
\chapter{Lógica de dominio}
% ============================================================================

\section{Máquina de estados de la persona}

Cada persona rastreada transita por los estados definidos en
\texttt{domain/constantes.py} y gestionados por \texttt{domain/persona.py}:

\begin{center}
\begin{tabular}{ll}
\toprule
\rowcolor{rvCodeBg}\textbf{Estado} & \textbf{Significado} \\
\midrule
En Tránsito & La persona circula, fuera de cualquier zona de interés. \\
Alcanzar Estante & Se acerca a una zona y aún no permanece el tiempo suficiente. \\
Inspeccionar Producto & Permanece quieta frente al estante (umbral de tiempo superado). \\
Mano en Estante & La muñeca entra al área de \textit{agarro} de la zona. \\
Producto Adquirido & Se confirma una toma (\textit{pick-up}). \\
Devolver el producto & Se confirma una devolución (\textit{put-back}) tras una toma previa. \\
\bottomrule
\end{tabular}
\end{center}

\section{Detección de pick-up / put-back}

La función \texttt{detectar\_interaccion\_manos} (\texttt{domain/persona.py}) implementa una
máquina de estados por \textbf{flancos}, no por umbrales instantáneos, para tolerar el
parpadeo del detector:

\begin{itemize}
  \item \textbf{Toma (pick-up):} se confirma cuando (a) la muñeca permanece
    \texttt{MANO\_EN\_ZONA\_PARA\_RETIRAR} \textit{frames} dentro del área de agarro \textbf{y}
    (b) el detector de objetos observa producto en posesión durante
    \texttt{OBJETO\_OBSERVACIONES\_CONFIRMAR} observaciones recientes. Existe un
    \textbf{respaldo por gesto} (\texttt{MANO\_EN\_ZONA\_CONFIRMA\_GESTO}) para productos que el
    modelo \textit{fine-tuneado} no reconoce.
  \item \textbf{Devolución (put-back):} exige una toma previa y, o bien (a) la muñeca
    \textbf{reingresa} al área de agarro (nueva visita) con el producto visto al entrar, o bien
    (b) el producto deja de observarse durante \texttt{OBJETO\_OBSERVACIONES\_REARMAR}
    observaciones mientras la mano sigue en el estante.
  \item \textbf{Deduplicación:} eventos del mismo tipo en la misma zona dentro de
    \texttt{VENTANA\_DEDUPE\_PICKUP} \textit{frames} se consideran el mismo evento real (el
    tracker puede fragmentar un ID en varios).
\end{itemize}

La exactitud de esta máquina de estados se mide con precisión/recall/F1 contra un
\textit{ground truth} anotado manualmente (\texttt{app/interfaces/validar\_embudo.py}).

\section{Clasificación de zonas y heatmap}

\texttt{domain/analisis.py} acumula, por zona, la \textbf{presencia ponderada por cercanía} al
centro de la zona y clasifica caliente / templada / fría por \textit{percentil} normalizado del
calor acumulado entre zonas. El heatmap de tráfico (\texttt{HeatmapTrafico}) usa un
\textit{splat} gaussiano local (no un recorrido completo del fotograma) por eficiencia, y
normaliza por frames con presencia (no por suma acumulada) para no saturar en videos largos.

\section{Agente de recomendación}

\texttt{domain/agente.py} es un motor \textbf{100\,\% determinista}: reglas sobre los datos
observados (visitas, tomas, tasa de rechazo, \textit{dwell time}). No usa modelos de lenguaje ni
llamadas externas. Genera texto accionable por zona y una propuesta de redistribución de
\textit{facing} (\% de espacio sugerido) cruzando la zona más caliente con la más fría.

% ============================================================================
\chapter{Privacidad y cumplimiento}
% ============================================================================

\section{Principio: privacidad por diseño}

El sistema cumple la \textbf{Ley 1581 de 2012} (Colombia, protección de datos personales)
mediante las siguientes garantías técnicas, verificadas por \texttt{scripts\_qa/qa\_privacidad.py}:

\begin{itemize}
  \item \textbf{Sin reconocimiento facial:} no se extraen \textit{embeddings} biométricos de
    ningún rostro.
  \item \textbf{Anonimización antes de cualquier render:} \texttt{infrastructure/anonimizacion.py}
    difumina (\textit{Gaussian blur}, kernel 61x61) la región de la cabeza, localizada con los
    \textit{keypoints} de YOLO-pose, \textbf{antes} de que el fotograma se muestre, se guarde o
    se transmita.
  \item \textbf{IDs efímeros:} los identificadores de tracking no re-identifican personas entre
    sesiones ni se cruzan con ninguna otra fuente de datos.
  \item \textbf{Procesamiento 100\,\% local (\textit{edge}):} ningún fotograma de video se
    envía a servicios en la nube de terceros para su análisis.
  \item \textbf{Salidas limpias:} el JSON de métricas y los fotogramas exportados no contienen
    información biométrica, solo coordenadas, tiempos y eventos agregados.
\end{itemize}

% ============================================================================
\chapter{Pruebas y calidad}
\label{sec:rendimiento}
% ============================================================================

\section{Suite de pruebas}

\begin{lstlisting}[language=bash]
python -m pytest tests/ -q                  # 51 tests de backend
python scripts_qa/qa_privacidad.py          # auditoria de privacidad (0 hallazgos)
python scripts_qa/qa_metricas.py            # invariantes de metricas
python scripts_qa/qa_rendimiento.py         # benchmark FPS / margen de tiempo real
cd frontend && npm test && npm run build    # tests y build del frontend
\end{lstlisting}

\begin{longtable}{>{\raggedright\arraybackslash}p{5.2cm} >{\raggedright\arraybackslash}p{8.4cm}}
\toprule
\rowcolor{rvCodeBg}\textbf{Módulo} & \textbf{Cobertura} \\
\midrule
\endhead
\texttt{tests/} (backend) & 51 pruebas: dominio (persona, polígonos, zonas), orquestador de
análisis (aislado, sin FastAPI ni subproceso real mediante ejecutores de prueba), API
(\texttt{TestClient}). \\
\texttt{frontend/tests/} & Pruebas de lógica de dominio del frontend (tasas, Sankey,
\textit{merge} de zonas) con \texttt{node:test}. \\
\texttt{qa\_privacidad.py} & Verifica ausencia de datos biométricos en las salidas. \\
\texttt{qa\_metricas.py} & Invariantes del embudo, rango del heatmap (0--100), coherencia de
recomendaciones. \\
\texttt{qa\_rendimiento.py} & Mide FPS de procesamiento y margen frente a tiempo real. \\
\bottomrule
\end{longtable}

\section{Resultados de referencia}

Medidos sobre GPU NVIDIA RTX 2060, video 720p: \textbf{\textasciitilde 24 FPS} de procesamiento,
equivalentes a un margen de \textbf{\textasciitilde 4$\times$} respecto al tiempo real del video
(es decir, un video de 1 minuto se procesa en \textasciitilde 15 segundos).

% ============================================================================
\chapter{Solución de problemas (técnico)}
% ============================================================================

\begin{longtable}{>{\raggedright\arraybackslash}p{5.2cm} >{\raggedright\arraybackslash}p{8.4cm}}
\toprule
\rowcolor{rvCodeBg}\textbf{Síntoma} & \textbf{Causa probable / solución} \\
\midrule
\endhead
\texttt{POST /api/analizar} responde 409 (\texttt{en\_curso}) & Ya hay un análisis en curso en el
mismo proceso; esperar a que termine o consultar \texttt{GET /api/estado}. \\
El video no genera vista previa & \texttt{cv2.VideoCapture} no pudo abrir el archivo (códec no
soportado); revisar \texttt{outputs/analysis.log}. \\
El pipeline termina con \texttt{estado=error} & Revisar \texttt{outputs/analysis.log}
(\texttt{SubprocessPipelineRunner} escribe ahí la salida completa del subproceso). \\
Las zonas no se reflejan en el análisis & Confirmar que \texttt{config/zonas\_runtime.json}
existe y es válido; \texttt{GET /api/zonas} expone el campo \texttt{fuente}. \\
Rendimiento bajo sin GPU & Esperado: el sistema funciona en CPU, pero a menor FPS. Ajustar
\texttt{PROCESAR\_CADA\_N} en \texttt{app/config.py} para procesar menos \textit{frames}. \\
\bottomrule
\end{longtable}

% ============================================================================
\chapter{Créditos y contacto}
% ============================================================================

\textbf{Docente asesor:} Sergio Arley Puerto Moreno

\textbf{Equipo de desarrollo:}

\begin{itemize}
  \item Silvana Sofía Siza Soriano
  \item Yeikov Jared Ramos Sánchez
  \item Sarita Toloza Buitrago
  \item Juan David Sandoval Bello
  \item Thomas Felipe Cepeda Mogollón
  \item Wendy Mariana Vargas Ramírez
  \item Andrés Felipe Galeano Téllez
  \item Sara Valentina Caicedo Duarte
\end{itemize}

\textbf{Institución:} Universidad Santo Tomás, Seccional Tunja --- Proyecto de Investigación /
Semillero RetailVision Edge AI.

\end{document}
